\documentclass[UTF8,a4paper,11pt,fontset=fandol]{ctexart}
\usepackage[left=72pt,right=72pt,top=78pt,bottom=78pt]{geometry}
\usepackage{amsmath,amsthm,booktabs,fancyhdr}
\setCJKmainfont[BoldFont=FandolHei-Bold.otf,ItalicFont=FandolKai-Regular.otf]{FandolHei-Regular.otf}
\renewcommand{\normalsize}{\fontsize{11pt}{17.05pt}\selectfont}
\ctexset{section={format=\fontsize{14pt}{20pt}\selectfont,beforeskip=12pt,afterskip=6pt}}
\pagestyle{fancy}
\fancyhf{}
\fancyhead[L]{设计语言论文样例}
\fancyfoot[C]{\thepage}
\setlength{\headheight}{14pt}
\newtheorem{定理}{定理}[section]
\numberwithin{equation}{section}
\title{设计作品中的结构化排版}
\author{豫言设计引擎}
\date{}
\begin{document}
\begin{center}
{\fontsize{18pt}{24pt}\selectfont 设计作品中的结构化排版}\par
\vspace{4pt}
豫言设计引擎
\end{center}
\noindent 摘要\par
本文讨论结构化文档与统一显示表的联系。文档先完成编号与引用求值，再按版心断行分页，最后由设计引擎统一输出。

\textbf{关键词：}设计语言；分页；结构化文档
\tableofcontents
\section{引言}\label{引言}
文章、界面与动画共享设计引擎。文档的引用关系在排版前求定，例如定理\ref{恒等}及公式\eqref{能量}。编号方法参见\cite{排版}。\footnote{本样例以点作为版式单位。}
\begin{定理}[恒等式]\label{恒等}
对于任意实数，恒等式成立。
\end{定理}
\begin{equation}\label{能量}
E=mc^2
\end{equation}
\begin{proof}
由符号定义直接得到所述等式。
\end{proof}
\clearpage
\section{分页与连续版面}\label{分页节}
分页版面按纸张高度转入下一页；连续版面重新流排，正文与编号保持同一来源。
\begin{table}[ht]
\centering
\caption{版面模式}\label{模式表}
\begin{tabular}{ll}\toprule
模式 & 用途 \\\midrule
分页 & 论文与印刷 \\
连续 & 网页与小屏 \\\bottomrule
\end{tabular}
\end{table}
\section*{索引}
分页\dotfill 2
\addcontentsline{toc}{section}{参考文献}
\begin{thebibliography}{1}
\bibitem{排版} 高德纳. 泰赫书[M]. 排版技术资料, 1984.
\end{thebibliography}
\end{document}
